\section{Radius-explicit stochastic first-driver bounds}
\label{S:radius-stochastic-anchor}

Retain every original covariance and amplitude in Hypothesis~\ref{H:homog},
with $a=\alpha_0$, and the real phase $E=H^{1/2}\times H^{-1/2}$.
Use the fixed root profile $\eta$ of Section~\ref{S:uniform-mean-regulation}
and physical mask $m_R=\eta(\cdot/R)^2$, $R\ge1$.
In the variance construction the base mask is $\zeta=\eta^2$.
The balanced unitary is
$J_0(u,v)=(\Lambda^{1/2}u,\Lambda^{-1/2}v)$.
Fix a finite horizon $B\ge1$ and parameters
\begin{equation}\label{E:radius-anchor-parameters}
  0<\gamma<\min\left(\frac12,\frac{2-\alpha}{2}\right),\qquad
  \frac12<\theta<\eta_t<\beta_t<
  \min\left(1-\frac a2,\frac{3-a-\alpha}{2}\right).
\end{equation}
The strict subcritical assumptions make these intervals nonempty.
Fix one even integer $Q\ge2$ such that
\begin{equation}\label{E:radius-anchor-Q}
  Q\gamma>d,\qquad
  \beta_t-\eta_t-\frac1Q>0,\qquad
  Q>\frac{2d}{1-a}.
\end{equation}
Such a $Q$ exists because all three lower thresholds are finite.
It stays fixed as $\delta$ tends to zero.
Constants below may depend on the model, $B,Q$, these exponents and the
fixed auxiliary cutoff profiles. They are independent of the spatial
radius, noise-basis truncation, moment index and deterministic residual
projection unless explicitly stated otherwise.

\subsection{Uniform outer factors and the exact whole-space envelope}

Fix real $\chi_0\in C_c^\infty(\R^d)$, equal to one on $B_1$ and
supported in the closed ball of radius two.
Choose real $\psi_0\in C_c^\infty(\R^d)$, equal to one on a
neighbourhood of $\operatorname{supp}\chi_0$ and supported inside $B_3$.
For $\rho\ge1$ put $\chi_\rho(x)=\chi_0(x/\rho)$,
$\psi_\rho(x)=\psi_0(x/\rho)$ and $L=2\pi\rho$.
Since $3<\pi$, both cutoffs are supported strictly inside the cube
of side $L$. The auxiliary torus has unnormalized $L^2(Q_L,\ud x)$,
orthonormal modes $L^{-d/2}e^{ik\cdot x/\rho}$ and Sobolev weights
$\langle k/\rho\rangle^r$.

\begin{lemma}[Uniform rescaled outer maps]\label{L:radius-outer-maps}
  The cutoff periodization and recovery maps from
  Lemma~\ref{L:cutoff-factorization} satisfy, for every real $r$,
  \[
    \|\mathcal P_{\psi_\rho}:H^r(\R^d)\to H^r(\mathbb T_L)\|
    \le C_{\psi_0,r},\qquad
    \|\mathcal E_{\chi_\rho}:H^r(\mathbb T_L)\to H^r(\R^d)\|
    \le C_{\chi_0,r},
  \]
  independently of $\rho\ge1$.
\end{lemma}

\begin{proof}
  The exact weighted Fourier kernels are
  \begin{align*}
    &L^{-d/2}
    \left(\frac{\langle k/\rho\rangle}{\langle\xi\rangle}\right)^r
    \widehat{\psi_\rho}(k/\rho-\xi),\\
    &L^{-d/2}
    \left(\frac{\langle\xi\rangle}{\langle k/\rho\rangle}\right)^r
    \widehat{\chi_\rho}(\xi-k/\rho).
  \end{align*}
  There is no additional $(2\pi)^{-d/2}$: it cancels between
  Fourier inversion and the ordinary integral in the first kernel,
  while the torus normalization gives the second.
  For either base profile $\phi_0$,
  $\widehat{\phi_\rho}(t)=\rho^d\widehat{\phi_0}(\rho t)$.
  Peetre's inequality~\eqref{E:cutoff-peetre}, also for negative $r$,
  and $\langle(\rho\xi-k)/\rho\rangle\le\langle\rho\xi-k\rangle$
  bound the relevant kernel by
  \[
    C_r\rho^{d/2}F_r(k-\rho\xi),\qquad
    F_r(t)=|\widehat{\phi_0}(\pm t)|\langle t\rangle^{|r|}.
  \]
  The nonnegative $F_r$ is integrable and bounded by arbitrarily
  high inverse powers of $\langle t\rangle$.
  At fixed $k$, integration in $\xi$ gives a bound $C\rho^{-d/2}$.
  At fixed $\xi$, summation in $k$ gives $C\rho^{d/2}$:
  translate $\rho\xi$ into a unit lattice cell and sum a power
  tail of order greater than $d$, as in~\eqref{E:cutoff-lattice-sum}.
  This lattice spacing is one and has no remaining $\rho$ dependence.
  The mixed Schur estimate bounds the operator norm by the square
  root of the product of those two quantities. Their radius powers
  cancel. For recovery the row and column roles are reversed.

  Bounded kernel extension from Schwartz functions and torus polynomials
  gives the full Sobolev maps by density.
  Their distributional meanings are the same cutoff maps as in
  Lemma~\ref{L:cutoff-factorization}; no point sampling of a rough
  Fourier function is used. This includes negative indices.
  This proves Lemma~\ref{L:radius-outer-maps}.
\end{proof}

\begin{proposition}[Radius-explicit compact envelopes]
  \label{P:radius-compact-envelopes}
  For real $\sigma$, $\gamma>0$ and $q\ge1$ with $q\gamma>d$,
  \begin{equation}\label{E:radius-cutoff-Schatten}
    \|\Lambda^\sigma\chi_\rho\Lambda^{-\sigma-\gamma}\|_{S_q}
    \le C_{\sigma,\gamma,q,d,\chi_0,\psi_0}\rho^{d/q}.
  \end{equation}
  In the balanced coordinates of~\eqref{E:compact-two-factors},
  \begin{equation}\label{E:radius-two-envelopes}
    \|K_{\rm in}\|_{S_q}=\|K_{\rm out}\|_{S_q}
    \le C_q\rho^{d/q}.
  \end{equation}
\end{proposition}

\begin{proof}
  The exact recovery identity, with $\chi_\rho\psi_\rho=\chi_\rho$,
  gives
  \[
    \Lambda^\sigma\chi_\rho\Lambda^{-\sigma-\gamma}
    =\Lambda^\sigma\mathcal E_{\chi_\rho}\mathcal I_\gamma
    \mathcal P_{\psi_\rho}\Lambda^{-\sigma-\gamma}.
  \]
  Here $\mathcal I_\gamma:H^{\sigma+\gamma}(\mathbb T_L)
  \to H^\sigma(\mathbb T_L)$ has diagonal singular values
  $\langle k/\rho\rangle^{-\gamma}$.
  The outer Bessel maps are unitary on their indicated Sobolev
  domains. The preceding lemma bounds both other outer maps
  uniformly in $\rho$.
  This is an equality on the full whole-space domain by the
  recovery proof and density, including negative $\sigma$.

  Put $u=q\gamma>d$. If $x\in k+[-1/2,1/2]^d$, then
  \[
    \langle x/\rho\rangle
    \le\sqrt2\langle k/\rho\rangle\langle(x-k)/\rho\rangle
    \le c_d'\langle k/\rho\rangle,\qquad
    c_d'=\sqrt{2(1+d/4)}.
  \]
  Integrate over the unit cubes to obtain
  \begin{equation}\label{E:radius-lattice-volume}
    \sum_{k\in\mathbb Z^d}\langle k/\rho\rangle^{-u}
    \le(c_d')^u\rho^d\int_{\R^d}\langle y\rangle^{-u}\ud y
    \le C_{d,u}\rho^d .
  \end{equation}
  The integral is finite: dyadic annuli outside the unit ball
  give a geometric series with terms bounded by $C_d2^{j(d-u)}$.

  Arrange the inclusion's diagonal values in decreasing order
  $(\lambda_j)$. They tend to zero. Retaining its first $j-1$
  coordinates has rank less than $j$ and error $\lambda_j$.
  Compose this approximation with the bounded outer factors.
  The singular-value approximation argument in
  Proposition~\ref{P:cutoff-singular-values} gives
  $s_j\le C_{\chi_0,\sigma}C_{\psi_0,\sigma+\gamma}\lambda_j$.
  Finite-rank approximation also proves compactness.
  Summing $q$th powers and using~\eqref{E:radius-lattice-volume}
  proves~\eqref{E:radius-cutoff-Schatten}.
  No separate Schatten ideal theorem is needed.

  The two blocks of $K_{\rm out}=P_{\chi_\rho}D_\gamma$
  are the whole-space operators with $\sigma=1/2,-1/2$.
  Their $q$th power norms add, proving its bound.
  Proposition~\ref{P:compact-two-envelopes} gives
  $K_{\rm in}=V K_{\rm out}V$ for the unitary coordinate swap
  $V(z_1,z_2)=(z_2,z_1)$, so their norms agree.
  In~\eqref{E:radius-anchor-parameters}, $\gamma<1/2$ and
  $q\gamma>d$ force $q>2d\ge2$; there is no Hilbert--Schmidt
  conclusion. Neither $P_{\chi_\rho}$ nor its adjoint is assumed
  selfadjoint or idempotent.
  This proves Proposition~\ref{P:radius-compact-envelopes}.
\end{proof}

\subsection{The actual driver and covariance contractions}

Write $M_R(W)=A_{m_R}(W)$ for the actual zero-start $E$ interaction
driver on $[0,B]$ from Proposition~\ref{P:actual-operator-driver}.
The phase remains the whole-space $E$ throughout.

\begin{proposition}[Radius-explicit actual path moments]
  \label{P:radius-actual-moments}
  Under~\eqref{E:radius-anchor-parameters}--\eqref{E:radius-anchor-Q},
  for every finite $v\ge1$,
  \begin{equation}\label{E:radius-actual-path-moments}
    \|M_R(W)\|_{L^v(C^\theta([0,B];\mathcal L(E)))}
    \le C_{Q,B}\sqrt v\,R^{d/Q},\qquad R\ge1.
  \end{equation}
  The constant does not depend on $v$ or $R$.
\end{proposition}

\begin{proof}
  Use the given raw smooth-test basis from
  Lemma~\ref{L:driver-noise-basis}, $g_j=W(f_j)$, and its
  actual finite paths $M_N=\sum_{j\le N}g_jA_j$.
  These are the same given-noise coefficient paths as in
  Lemma~\ref{L:driver-holder-limit}, not source approximations.
  Choose
  \[
    \rho=2R+B+2\le(B+4)R,\qquad \chi=\chi_0(\cdot/\rho).
  \]
  The plateau covers $B_{2R+B+1}$.
  Thus both literal localization identities
  \eqref{E:driver-actual-localization} hold for this envelope.
  The physical noise mask remains $m_R$.

  For an interval $I\subset[0,B]$ of length $0<h\le1$,
  Proposition~\ref{P:wave-two-variances} and
  Lemma~\ref{L:driver-actual-variances} give both finite weighted
  variances with budget $C_Bh^{2-a}B_\alpha(h)$,
  uniformly in $R,N,I$.
  In particular, the base mask is the same fixed $\eta^2$
  required by that variance theorem. Corollary~\ref{C:finite-gaussian-operators},
  applied in the balanced coordinates, gives for every even $q\ge Q$
  \[
    \|M_N(I)\|_{L^q(S_q)}
    \le\sqrt{q-1}\sqrt{C_B}\,
    h^{1-a/2}B_\alpha(h)^{1/2}\|K_{\rm out}\|_{S_q},
  \]
  where $B_\alpha(h)$ is $1$, $1+\log(1/h)$ or $h^{1-\alpha}$
  according as $\alpha<1$, $\alpha=1$ or $\alpha>1$.
  Both actual variances and both localizations have been checked;
  individual basis-dependent coefficient bounds are not summed.

  For $q\ge Q$, the unnormalized Schatten norms obey
  $\|K_{\rm out}\|_{S_q}\le\|K_{\rm out}\|_{S_Q}$.
  Indeed, normalize a nonzero singular-value sequence by its
  $\ell^Q$ norm; each entry is at most one, so its $q$th power
  sum is at most its $Q$th power sum. The zero case is immediate.
  Apply~\eqref{E:radius-two-envelopes} only at the fixed index $Q$.
  Since $\beta_t$ is strictly below the corresponding short-time
  power, it follows that
  \begin{equation}\label{E:radius-finite-increment}
    \|M_N(t)-M_N(u)\|_{L^q(\mathcal L(E))}
    \le C_{Q,B}\sqrt q\,R^{d/Q}|t-u|^{\beta_t},
    \qquad |t-u|\le1.
  \end{equation}
  At $\alpha=1$, the factor $\sqrt{1+\log(1/h)}$ is absorbed
  by the strict gap $1-a/2-\beta_t>0$.
  For $\alpha>1$, the power is $(3-a-\alpha)/2$.
  These power/log constants are independent of $R,N,q$.
  Longer increments are finite sums of increments of length at
  most one, with constants depending only on the fixed $B$.

  Let $D_n(F)$ be the largest adjacent increment on the grid
  $kB2^{-n}$ of $[0,B]$.
  For fine levels $B2^{-n}\le1$, the finite-maximum bound gives
  \[
    \|D_n(M_N)\|_{L^q}
    \le2^{n/q}C_{Q,B}\sqrt q\,R^{d/Q}(B2^{-n})^{\beta_t}.
  \]
  No independence of increments is used. The sum after multiplication
  by $2^{n\eta_t}$ converges uniformly for even $q\ge Q$, because
  its geometric ratio is at most
  $2^{-(\beta_t-\eta_t-1/Q)}<1$.
  The finitely many coarse levels are controlled by unit partitions.
  The deterministic dyadic chaining argument in
  Lemma~\ref{L:driver-holder-limit} therefore bounds the
  $L^qC^{\eta_t}$ norms by $C_{Q,B}\sqrt q\,R^{d/Q}$.
  Each finite path is continuous and starts at zero, as required there.
  The fixed-horizon embedding $C^{\eta_t}\hookrightarrow C^\theta$
  gives the same estimate in $C^\theta$.

  For each fixed $R$, Lemma~\ref{L:driver-holder-limit} identifies
  $M_R(W)$ as the all-finite-moment $C^\theta$ limit of these
  very $M_N$. Pass the uniform finite-$N$ estimate to this limit.
  Uniform convergence in $R$ is unnecessary: the limiting bound has
  its explicit radius dependence at every $R$.
  For arbitrary $v\ge1$, choose an even $q\ge\max(v,Q)$ with
  $q\le v+Q+2\le(Q+3)v$.
  Lyapunov's inequality gives~\eqref{E:radius-actual-path-moments},
  absorbing $\sqrt{Q+3}$ into the fixed constant.
  This proves Proposition~\ref{P:radius-actual-moments}.
\end{proof}

\begin{corollary}[Actual deterministic covariance residuals]
  \label{C:radius-covariance-residuals}
  Let $Y$ be a centered Gaussian linear noise on the actual coupled
  space, with $\operatorname{Var}Y(f)\le C\|f\|_{H_W}^2$, $C\ge0$.
  Its actual transferred driver satisfies
  \[
    \|M_R(Y)\|_{L^v(C^\theta)}
    \le\sqrt C\,C_{Q,B}\sqrt v\,R^{d/Q}.
  \]
  In particular this holds with $C=1$ for
  $Y_m(f)=W((I-\Pi)T_{\delta,m}f)$, for any deterministic
  orthogonal projection $\Pi$ on the slow $H_W$.
  Its rank may be infinite and it may depend on $\delta,R$.
\end{corollary}

\begin{proof}
  Proposition~\ref{P:source-gaussian-transfer} is a linear Gaussian
  covariance contraction theorem for this same physical mask and
  these same deterministic coefficient maps.
  It bounds the path norm by $\sqrt C$ times the reference
  isonormal driver norm. The reference has the canonical first-driver
  law used in Proposition~\ref{P:radius-actual-moments}; the actual
  $Y$ is not replaced in this comparison. The zero case gives zero
  directly. This proves the first assertion.
  For the second, the actual raw isometry
  \eqref{E:raw-block-pullback} and $\|I-\Pi\|\le1$ give
  $\operatorname{Var}Y_m(f)\le\|f\|_{H_W}^2$,
  where the right side is the original raw norm in fast coordinates.
  The joint residuals retain the given slow $W$ and may be correlated.
  They need not lie in the span of an individual fast noise,
  which is not a hypothesis of Gaussian transfer.
  No random noise-dependent projection is used.
  This proves Corollary~\ref{C:radius-covariance-residuals}.
\end{proof}

\subsection{The growing-radius expected-norm anchor}

\begin{proposition}[Actual averaged block mean tends to zero]
  \label{P:radius-mean-anchor}
  Use the true block partition $B,\delta,N,\ell_m$ of
  Section~\ref{S:actual-cameron-mean}, with
  $0<\delta<\min(1,1/B)$ and $N=\lfloor1/(\delta B)\rfloor$.
  For $Y_m$ as in Corollary~\ref{C:radius-covariance-residuals},
  put
  \[
    X_m(u)=g_\delta M_R(Y_m)(\min(u,\ell_m)),\qquad
    Z_\delta=\left(\delta\sum_{m=0}^N
    \|X_m\|_{C^\theta}^p\right)^{1/p},\qquad
    g_\delta=\delta^{(1-a)/2},
  \]
  for any fixed $p\ge1$. Then
  \begin{equation}\label{E:radius-mean-anchor}
    \E Z_\delta\le(2/B)^{1/p}C_{Q,B}\sqrt p\,
    g_\delta R^{d/Q}.
  \end{equation}
  In particular, for $R_\delta=\delta^{-1}+2$,
  \begin{equation}\label{E:radius-mean-anchor-decay}
    2\E Z_\delta
    \le2(2/B)^{1/p}C_{Q,B}\sqrt p\,3^{d/Q}
    \delta^{e_{\rm anchor}}\longrightarrow0,\qquad
    e_{\rm anchor}=\frac{1-a}{2}-\frac dQ>0 .
  \end{equation}
  Both bounds are uniform in the deterministic orthogonal projection.
\end{proposition}

\begin{proof}
  Stopping a zero-start interaction primitive and extending it
  constantly to $B$ does not increase its supremum or H\"older
  seminorm. An empty last interval gives the zero path.
  Probability normalization and Lyapunov give
  \begin{align*}
    \E Z_\delta
    &\le(\E Z_\delta^p)^{1/p}
    =\left(\delta\sum_{m=0}^N
    \E\|X_m\|_{C^\theta}^p\right)^{1/p}\\
    &\le[\delta(N+1)]^{1/p}g_\delta
    C_{Q,B}\sqrt p\,R^{d/Q}.
  \end{align*}
  This is linearity for a finite sum; no time-block independence
  or maximum over blocks is used.
  Since $\delta B<1$, we have $N\ge1$ and
  $N+1\le2/(\delta B)$, proving~\eqref{E:radius-mean-anchor}.

  Now $R_\delta\le3/\delta$ because $\delta<1$.
  The last condition in~\eqref{E:radius-anchor-Q} is equivalent to
  $Q(1-a)>2d$. Dividing by $2Q>0$ gives
  $e_{\rm anchor}>0$.
  Substitute this radius into~\eqref{E:radius-mean-anchor}
  to obtain~\eqref{E:radius-mean-anchor-decay}.
  The estimate holds at each fixed $B,p,Q$ and model, uniformly
  in every deterministic residual projection.
  This proves Proposition~\ref{P:radius-mean-anchor}.
\end{proof}

These bounds concern the expected first-driver norm on $E$.
They do not give an $F$ estimate with constant initial-data columns.
The radius substitution is numerical; cone coverage and source-to-point
transfer require their separate arguments.
No concentration, infinite-head Gaussian integration or nonlinear moment
upper bound follows from this section alone.
The constants need not be uniform as $Q,B$ tend to infinity or as model
parameters approach an endpoint. For zero covariance, the actual drivers
and stochastic bounds vanish directly.
